\chapter{React State Management}

\section{The Core Hooks in Practice}

\begin{lstlisting}[language=JavaScript]
// useState - local, per-render state
const [tasks, setTasks] = useState([]);

// useEffect - side effects (fetching, subscriptions, timers)
useEffect(() => {
  fetchTasks();
}, []); // empty array = run once, on mount

// useContext - read shared state without prop drilling
const { user, logout } = useContext(AuthContext);

// useRef - mutable value that survives re-renders, no re-render on change
const inputRef = useRef(null);
useEffect(() => { inputRef.current?.focus(); }, []);

// useMemo - cache an expensive computed value
const doneCount = useMemo(
  () => tasks.filter((t) => t.status === "done").length,
  [tasks]
);

// useCallback - cache a function identity across re-renders
const handleDelete = useCallback(
  (id) => setTasks((prev) => prev.filter((t) => t._id !== id)),
  []
);
\end{lstlisting}

\section{A Custom Hook: useAuth}

\begin{lstlisting}[language=JavaScript]
// hooks/useAuth.js
import { useContext } from "react";
import { AuthContext } from "../context/AuthContext";

export function useAuth() {
  const context = useContext(AuthContext);
  if (!context) {
    throw new Error("useAuth must be used within an AuthProvider");
  }
  return context; // { user, loading, login, logout }
}
\end{lstlisting}

Custom hooks let you extract stateful logic (auth status, form handling,
debounced search) into a reusable function without duplicating
\texttt{useEffect}/\texttt{useState} wiring across components.

\section{Common Mistakes: Infinite API Call Loops}

The single most common MERN bug is an effect that re-triggers itself forever.

\begin{warnbox}[WARNING: Infinite Loop]
\begin{lstlisting}[language=JavaScript]
// BUG: no dependency array -> runs after every render,
// and setTasks triggers another render -> infinite loop
function Dashboard() {
  const [tasks, setTasks] = useState([]);

  useEffect(() => {
    getTasks().then(setTasks);
  }); // <-- missing [] !!

  return <TaskList tasks={tasks} />;
}
\end{lstlisting}
Every call to \texttt{setTasks} causes a re-render, which runs the effect again
(no dependency array means ``run after every render''), which calls
\texttt{setTasks} again -- forever. This shows up as a browser tab hammering
your API with hundreds of requests per second.
\end{warnbox}

\textbf{Fix:} add the dependency array, even if empty.

\begin{lstlisting}[language=JavaScript]
useEffect(() => {
  getTasks().then(setTasks);
}, []); // runs once, on mount
\end{lstlisting}

A second, subtler version of the same bug happens with objects or functions
in the dependency array:

\begin{warnbox}[WARNING: Unstable Dependency]
\begin{lstlisting}[language=JavaScript]
// BUG: filters is a new object every render -> effect
// "sees" a new dependency every time -> runs every render
function TaskList() {
  const filters = { status: "pending" }; // new reference each render

  useEffect(() => {
    getTasks(filters).then(setTasks);
  }, [filters]); // filters is never === the previous filters
}
\end{lstlisting}
\end{warnbox}

\textbf{Fix:} memoize the object, or depend on its primitive fields instead.

\begin{lstlisting}[language=JavaScript]
const filters = useMemo(() => ({ status }), [status]);

useEffect(() => {
  getTasks(filters).then(setTasks);
}, [filters]);

// or simply:
useEffect(() => {
  getTasks({ status }).then(setTasks);
}, [status]); // depend on the primitive, not the wrapper object
\end{lstlisting}

\begin{tipbox}[TIP]
Rule of thumb: every value read inside \texttt{useEffect} that can change over
time belongs in the dependency array. If adding it causes a loop, the real
fix is almost always to memoize that value or depend on a primitive
derived from it -- not to remove it from the array.
\end{tipbox}
